\section{System Architecture \& Data Flow}
\label{sec:system_architecture}

\subsection{Architecture Overview}
\label{subsec:architecture_overview}

GridSense AI is structured around a modular, multi-tier architecture designed to decouple data ingestion, machine learning forecasting, risk intelligence, mathematical optimization, and operator interaction. 

Rather than deploying a monolithic prediction service, the system isolates each analytical stage into dedicated computational engines connected through a shared data access layer. \Cref{fig:system_architecture} illustrates the high-level system architecture.

\begin{figure}[htbp]
    \centering
    \begin{tikzpicture}[
        node distance=0.42cm and 0.4cm,
        layerBox/.style={rectangle, rounded corners=3pt, draw=accentTeal!80, fill=accentTeal!10, text width=8.8cm, align=center, font=\footnotesize, inner sep=4pt},
        dataBox/.style={rectangle, rounded corners=3pt, draw=accentBlue!80, fill=accentBlue!10, text width=2.65cm, align=center, font=\scriptsize, inner sep=3.5pt},
        dbBox/.style={rectangle, rounded corners=3pt, draw=primaryNavy!80, fill=boxBg, text width=2.2cm, align=center, font=\scriptsize, inner sep=4pt},
        dashBox/.style={rectangle, rounded corners=3pt, draw=accentAmber!80, fill=accentAmber!10, text width=8.8cm, align=center, font=\footnotesize, inner sep=4.5pt},
        arrow/.style={-{Stealth[scale=0.85]}, thick, color=darkSlate},
        darrow/.style={{Stealth[scale=0.85]}-{Stealth[scale=0.85]}, dashed, thick, color=primaryNavy}
    ]
        % Top Inputs
        \node[dataBox] (src1) {\textbf{Weather Data}\\Irradiance, Wind, Temp, Cloud Cover};
        \node[dataBox, right=0.35cm of src1] (src2) {\textbf{Historical Records}\\Solar/Wind Generation Time Series};
        \node[dataBox, right=0.35cm of src2] (src3) {\textbf{Site Parameters}\\Installed Capacity, Coordinates, Inverter Spec};

        % Ingestion Layer
        \node[layerBox, below=0.55cm of src2] (ingest) {\textbf{Data Ingestion \& Validation Layer}\\Format Checking \textbar{} Missing Value Imputation \textbar{} Timestamp Alignment};

        % Forecasting Layer
        \node[layerBox, below=of ingest] (forecast) {\textbf{Forecasting Engine}\\Multi-Horizon Generation Prediction (24h, 48h, 72h)};

        % Uncertainty & Risk Layer
        \node[layerBox, below=of forecast] (risk) {\textbf{Uncertainty \& Risk Engine}\\Quantile Intervals ($p_{10}, p_{50}, p_{90}$) \textbar{} Surplus/Shortfall Identification};

        % Simulation & Optimization Layer
        \node[layerBox, below=of risk] (optim) {\textbf{Simulation \& Decision Optimization Engine}\\What-If Scenario Stress Testing \textbar{} Storage Dispatch \& Curtailment Optimization};

        % Recommendation Layer
        \node[layerBox, below=of optim] (recom) {\textbf{Recommendation Layer \& Explainable Copilot}\\Actionable Operational Schedules \textbar{} Cost \& Carbon Contextualization};

        % Operator Dashboard
        \node[dashBox, below=of recom] (dash) {\textbf{Operator Dashboard \& Interaction Layer}\\Multi-Horizon Visualizations \textbar{} Risk Timeline \textbar{} Human Decision Approval};

        % Central Data Store (Right Side)
        \node[dbBox, right=0.5cm of risk] (db) {\textbf{Central Data\\\& State Store}\\[0.2em]\tiny Raw Logs\\Historical Store\\Forecast Cache\\Scenario State};

        % Arrows
        \draw[arrow] (src1) -- (src1 |- ingest.north);
        \draw[arrow] (src2) -- (ingest.north);
        \draw[arrow] (src3) -- (src3 |- ingest.north);

        \draw[arrow] (ingest) -- (forecast);
        \draw[arrow] (forecast) -- (risk);
        \draw[arrow] (risk) -- (optim);
        \draw[arrow] (optim) -- (recom);
        \draw[arrow] (recom) -- (dash);

        % DB interactions
        \draw[darrow] (ingest.east) -| (db.north);
        \draw[darrow] (forecast.east) -- (db.west |- forecast.east);
        \draw[darrow] (risk.east) -- (db.west);
        \draw[darrow] (optim.east) -- (db.west |- optim.east);
        \draw[darrow] (recom.east) -| (db.south);
    \end{tikzpicture}
    \caption{GridSense AI Layered System Architecture and Central Data Interaction.}
    \label{fig:system_architecture}
\end{figure}

The architecture guarantees that raw data transformations, predictive modeling, operational risk scoring, and optimization routines remain independently testable, modular, and maintainable.

\subsection{Data Sources}
\label{subsec:data_sources}

The platform ingests three primary data categories explicitly specified in the HackOut'26 problem statement:
\begin{itemize}[leftmargin=1.8em, itemsep=0.25em]
    \item \textbf{Weather Data:} High-resolution numerical weather predictions and environmental feeds, including global horizontal irradiance (GHI), direct normal irradiance (DNI), ambient temperature, cloud cover percentages, wind speed at hub height, and barometric pressure.
    \item \textbf{Historical Generation Data:} Timestamped time-series records of past solar irradiance power output and wind turbine power generation, capturing diurnal patterns, seasonal baselines, and historical ramps.
    \item \textbf{Site-Level Parameters:} Metadata defining the physical installation, including nominal installed capacity (MW), geographic coordinates, solar panel tilt and azimuth, wind turbine hub heights, power curves, and inverter efficiency ratings.
\end{itemize}

In addition to these mandatory inputs, GridSense AI accommodates \textbf{operational parameters} (such as battery energy storage capacity, current state of charge, grid interconnection limits, and standby backup generator ramp rates) as supplementary inputs to the decision optimization layer.

\subsection{Data Ingestion \& Validation}
\label{subsec:data_ingestion}

Incoming data streams are ingested through automated validation pipelines before being passed to downstream analytical models. The ingestion layer executes four critical operations:
\begin{enumerate}[leftmargin=1.8em, itemsep=0.2em]
    \item \textbf{Schema and Format Validation:} Verifying data structures, variable names, and expected data types.
    \item \textbf{Missing-Value Imputation:} Detecting telemetry drops, sensor dropouts, or communication interruptions and applying localized time-series interpolations.
    \item \textbf{Timestamp Alignment \& Resampling:} Synchronizing disparate weather forecasts and power telemetry into uniform hourly intervals across UTC and local grid time.
    \item \textbf{Unit and Boundary Consistency:} Normalizing units (e.g., converting $\text{W/m}^2$, $\text{m/s}$, and $\text{kW}$ to standardized $\text{MW}$ indices) and flagging non-physical outliers (such as negative solar irradiance at night).
\end{enumerate}

Isolating validation at ingestion ensures downstream forecasting models receive clean, synchronized, and physically consistent feature matrices.

\subsection{Forecasting Layer}
\label{subsec:forecasting_layer}

The forecasting layer combines validated weather features, historical generation baselines, and temporal characteristics (hour of day, day of year) to project renewable power generation across 24-hour, 48-hour, and 72-hour planning horizons.

Importantly, the engine does not produce solitary point estimates in isolation. It computes probabilistic prediction intervals ($p_{10}$, $p_{50}$, $p_{90}$) alongside point forecasts, quantifying the statistical envelope of potential generation outcomes:
\begin{align*}
    \text{Expected Generation } (p_{50}) &= 80\text{ MW} \\
    \text{Prediction Interval } [p_{10}, p_{90}] &= [72\text{ MW}, 88\text{ MW}] \\
    \text{Assessed Uncertainty Level} &= \text{Moderate}
\end{align*}

These uncertainty bands form the empirical foundation for downstream risk detection.

\subsection{Intelligence \& Decision Layer}
\label{subsec:intelligence_decision}

The intelligence and decision layer ingests generation forecasts and uncertainty envelopes to identify potential operational grid events. Key detected conditions include:
\begin{itemize}[leftmargin=1.8em, itemsep=0.2em]
    \item \textbf{Renewable Shortfall Events:} Projected deficit periods where renewable output falls below baseline demand.
    \item \textbf{Renewable Surplus Events:} Anticipated over-generation periods where clean generation exceeds local absorption capacity.
    \item \textbf{High-Volatility Ramps:} Time intervals characterized by steep generation deltas or wide prediction intervals.
    \item \textbf{Curtailment and Storage Windows:} Identification of optimal hours for battery charging versus risk of forced renewable spillage.
\end{itemize}

Upon event detection, the decision engine executes multi-variable optimization to evaluate candidate operational responses (storage dispatch, power export, backup scheduling, or managed curtailment), balancing operational expense against renewable utilization.

\subsection{User Interaction Layer}
\label{subsec:user_interaction}

The operator dashboard and AI Copilot deliver a unified, interactive decision-support interface. The user interaction layer presents:
\begin{itemize}[leftmargin=1.8em, itemsep=0.2em]
    \item \textbf{Multi-Horizon Forecast Curves:} Interactive generation timelines with overlaid demand baselines and confidence bounds.
    \item \textbf{Operational Risk Timeline:} Color-coded alert banners pinpointing upcoming surplus, shortfall, and high-ramp periods.
    \item \textbf{Recommended Action Cards:} Clear, actionable schedules detailing recommended storage dispatch, backup commitment, and curtailment setpoints.
    \item \textbf{Explainable AI Copilot:} Natural-language summaries explaining the meteorological drivers, risk factors, and financial/carbon impacts of recommended actions.
    \item \textbf{What-If Scenario Sandbox:} Interactive controls allowing operators to perturb weather variables or asset availability and inspect updated recommendations in real time.
\end{itemize}

\subsection{End-to-End Data Flow}
\label{subsec:end_to_end_flow}

The lifecycle of information through GridSense AI is organized into eight synchronized sequential steps, as illustrated in \cref{fig:data_flow_steps}:
\begin{enumerate}[leftmargin=1.8em, itemsep=0.2em]
    \item \textbf{Data Collection:} Continuous ingestion of meteorological forecasts, plant telemetry, and site parameters.
    \item \textbf{Data Validation:} Automated cleansing, outlier removal, timestamp alignment, and normalization.
    \item \textbf{Forecast Generation:} Multi-horizon generation forecasting across 24h, 48h, and 72h windows.
    \item \textbf{Uncertainty \& Risk Analysis:} Computation of confidence intervals, risk scoring, and surplus/shortfall detection.
    \item \textbf{Scenario Simulation:} Dynamic evaluation of operator-defined what-if perturbations and edge cases.
    \item \textbf{Decision Evaluation:} Multi-objective optimization of storage, backup, import/export, and curtailment actions.
    \item \textbf{Explainable Recommendation:} Generation of natural-language operational guidance with cost and carbon trade-offs.
    \item \textbf{Operator Decision Support:} Presentation to qualified control-room operators for human review, verification, and dispatch approval.
\end{enumerate}

\begin{figure}[htbp]
    \centering
    \begin{tikzpicture}[
        node distance=0.35cm,
        stepBox/.style={rectangle, rounded corners=3pt, draw=accentTeal!80, fill=boxBg, text width=7.4cm, align=left, font=\scriptsize, inner sep=3.5pt},
        numBox/.style={rectangle, rounded corners=2pt, draw=accentTeal, fill=accentTeal, text=white, font=\scriptsize\bfseries, inner sep=2pt},
        arrow/.style={-{Stealth[scale=0.8]}, thick, color=darkSlate}
    ]
        \node[stepBox] (s1) {\textbf{1. Data Collection} \hfill \color{darkSlate!70}Weather, Historical Generation, Site Parameters};
        \node[stepBox, below=of s1] (s2) {\textbf{2. Data Validation} \hfill \color{darkSlate!70}Format Check, Imputation, Unit Alignment};
        \node[stepBox, below=of s2] (s3) {\textbf{3. Forecast Generation} \hfill \color{darkSlate!70}24h, 48h, 72h Multi-Horizon Predictions};
        \node[stepBox, below=of s3] (s4) {\textbf{4. Uncertainty \& Risk Analysis} \hfill \color{darkSlate!70}Confidence Intervals, Surplus/Deficit Windows};
        \node[stepBox, below=of s4] (s5) {\textbf{5. Scenario Simulation} \hfill \color{darkSlate!70}What-If Stress Testing, Parameter Adjustments};
        \node[stepBox, below=of s5] (s6) {\textbf{6. Decision Evaluation} \hfill \color{darkSlate!70}Storage, Backup, Import/Export, Curtailment};
        \node[stepBox, below=of s6] (s7) {\textbf{7. Action Recommendation} \hfill \color{darkSlate!70}Explainable Guidance, Cost \& Carbon Insights};
        \node[stepBox, below=of s7, draw=accentAmber!90, fill=accentAmber!10] (s8) {\textbf{8. Operator Decision Support} \hfill \color{darkSlate!70}\textbf{Human Review, Approval \& Dispatch}};

        \draw[arrow] (s1) -- (s2);
        \draw[arrow] (s2) -- (s3);
        \draw[arrow] (s3) -- (s4);
        \draw[arrow] (s4) -- (s5);
        \draw[arrow] (s5) -- (s6);
        \draw[arrow] (s6) -- (s7);
        \draw[arrow] (s7) -- (s8);
    \end{tikzpicture}
    \caption{Operational End-to-End Data Flow through the GridSense AI Platform.}
    \label{fig:data_flow_steps}
\end{figure}

\subsection{Phase 3 Conclusion}
\label{subsec:phase3_conclusion}

GridSense AI employs a modular architecture where each component maintains a distinct, well-defined operational responsibility. 

Data ingestion establishes a clean empirical foundation; the forecasting layer generates multi-horizon predictions with quantified uncertainty; the intelligence layer identifies imminent grid imbalances; the optimization engine evaluates multi-criteria dispatch strategies; and the dashboard empowers the operator with explainable, actionable recommendations.

This architectural framework provides the foundation for the detailed AI/ML forecasting pipelines and decision models examined in the subsequent phases.
